\section{Limits and Scope Boundaries}

This section specifies the boundaries of applicability of the framework developed in this paper and clarifies its explicit non-goals. The purpose is to delineate what the analysis does and does not claim, thereby preventing overextension of the results beyond their intended scope.

First, the framework addresses STOP exclusively as a diagnostic outcome at the level of an enterprise model. It does not claim to describe, predict, or evaluate organizational decisions, behaviors, or outcomes in practice. Any correspondence between a STOP diagnosis and real-world action lies outside the scope of the analysis. This restriction follows from the distinction between model-level representability and real-world execution that characterizes formal approaches to anticipatory systems \cite{rosen1985}.

Second, the analysis is model-relative by construction. All claims concerning STOP depend on the specification of an enterprise model and its internal constraints. The framework does not provide criteria for selecting, validating, or optimizing such models, nor does it assert that any particular model is correct or complete. Questions of model choice and model adequacy are explicitly excluded.

Third, the framework does not address the process by which models are revised, extended, or replaced following a STOP outcome. While STOP may imply that continuation would require a change of model, the nature, justification, or governance of such changes is not considered here.

Fourth, no claims are made regarding the frequency, likelihood, or empirical distribution of STOP outcomes in real enterprises. The framework establishes logical possibility and formal validity, not empirical prevalence or statistical expectation.

Finally, the analysis does not attempt to resolve normative questions. It does not state whether STOP outcomes are desirable or undesirable, nor does it evaluate their organizational, economic, or social consequences. Such evaluations require additional criteria and perspectives not introduced in this work and therefore lie beyond the scope of inquiry addressed here \cite{churchman1971}.

Within these boundaries, the framework establishes a precise and limited result: STOP can be a formally valid diagnostic outcome when enterprises are modeled as continua and admissibility of continuation is exhausted. All interpretations beyond this claim fall outside the intended scope of the paper.
